\documentclass[12pt,a4paper]{report}
\usepackage[margin=25mm]{geometry}
\usepackage{graphicx}
\usepackage{float}
\usepackage{amsmath,amssymb}
\usepackage{booktabs}
\usepackage{array}
\usepackage{hyperref}
\usepackage{xcolor}
\usepackage{enumitem}
\usepackage{caption}

\graphicspath{{report_figures/}}
\hypersetup{colorlinks=true, linkcolor=blue, citecolor=blue, urlcolor=blue}

\newcommand{\writingbox}[1]{
  \vspace{0.5em}
  \noindent\fbox{\parbox{0.96\linewidth}{\textbf{Writing space:} #1\\[6em]}}
  \vspace{0.8em}
}

\newcommand{\figuretodo}[3]{
  \begin{figure}[H]
    \centering
    \includegraphics[width=#1\linewidth]{#2}
    \caption{#3}
  \end{figure}
}

\begin{document}

\begin{titlepage}
  \centering
  \vspace*{2cm}
  {\Large University of Cambridge Department of Engineering\par}
  \vspace{1.5cm}
  {\Huge\bfseries Neuromorphic Bat-Inspired 3D Echolocation for Edge Localisation\par}
  \vspace{1.5cm}
  {\Large Project Draft Plan\par}
  \vfill
  \begin{tabular}{ll}
    Name: & \textit{[Your name]}\\
    College: & \textit{[Your college]}\\
    Supervisor: & \textit{[Supervisor]}\\
    Date: & \today
  \end{tabular}
\end{titlepage}

\pagenumbering{roman}
\tableofcontents
\clearpage

\chapter*{Technical Abstract Plan}
\addcontentsline{toc}{chapter}{Technical Abstract Plan}
\textbf{Required length:} no more than two pages. Include name, College, and project title because this is submitted separately.

\writingbox{Write this last. Summarise the problem, the biomimetic approach, the neuromorphic primitives, the three parallel pathways, the CANN readout, headline results, and limitations/future work.}

\chapter*{Chapter 0: Acronyms}
\addcontentsline{toc}{chapter}{Chapter 0: Acronyms}
\begin{tabular}{ll}
  AC & Auditory cortex\\
  CANN & Continuous attractor neural network\\
  CD & Corollary discharge\\
  COM & Centre of mass\\
  DCN & Dorsal cochlear nucleus\\
  DNLL & Dorsal nucleus of the lateral lemniscus\\
  FFT & Fast Fourier transform\\
  IC & Inferior colliculus\\
  IIR & Infinite impulse response\\
  ILD & Interaural level difference\\
  ITD & Interaural time difference\\
  LIF & Leaky integrate-and-fire\\
  LSO & Lateral superior olive\\
  MNTB & Medial nucleus of the trapezoid body\\
  PSD & Power spectral density\\
  RF & Resonate-and-fire\\
  SC & Superior colliculus\\
  SNN & Spiking neural network\\
  VCN & Ventral cochlear nucleus
\end{tabular}

\pagenumbering{arabic}

\chapter{Introduction and Motivation}
\section{The Engineering Problem}
\writingbox{Describe 3D sound localisation as estimating range, azimuth, and elevation from echoes under edge-system constraints. Explain why low latency, low power, and robustness matter.}

\figuretodo{0.82}{ch1_emitted_call_spectrogram.png}{Suggested introductory figure: emitted frequency-modulated call spectrogram.}
\figuretodo{0.82}{ch1_binaural_head_shadow.png}{Suggested figure: binaural received signals and head-shadow cue.}

\section{Radar/Sonar Context}
\writingbox{Explain classical ranging, matched filtering, beamforming, ITD/ILD analogues, and where conventional approaches become expensive.}

\section{Biomimetic Pivot}
\writingbox{Explain why bat echolocation is the biological reference system: small power budget, real-time performance, frequency sweeps, parallel auditory pathways, and spike-based computation.}

\section{Architectural Philosophy}
\writingbox{State the report's design thesis: bottom-up functional tuning first, then targeted training if needed. Contrast this with opaque end-to-end training.}

\chapter{Neuromorphic Primitives and Classical Algorithmic Parallels}
\section{Spikes and Neuron Models}
\figuretodo{0.8}{ch2_lif_microdynamics.png}{LIF microdynamics.}
\figuretodo{0.8}{ch2_rf_microdynamics.png}{RF microdynamics.}
\figuretodo{0.8}{ch2_frequency_gain.png}{LIF, RF, and level-crossing frequency response comparison.}
\writingbox{Explain LIF, RF, and level-crossing neurons. Include equations and functional interpretation.}

\section{Neurons as Algorithms}
\writingbox{Map RF banks to band-pass/Fourier-like analysis, LIF coincidence banks to cross-correlation, synaptic weights to matched filters, and population spikes to probabilistic samples.}

\section{Preliminary Coincidence-Detection Control}
\figuretodo{0.82}{ch2_sustained_pitch_response.png}{Suggested control result: sustained-pitch response of LIF/RF/binary coincidence detectors.}
\writingbox{Present the LIF/RF/binary distance coincidence tests. Use this to justify why LIF/binary coincidence is used for precise timing and RF is not used as the primary range detector.}

\chapter{System Architecture and Peripheral Frontend}
\section{End-to-End Pipeline}
\writingbox{Introduce the full architecture: signal generation, echo simulation, cochlea, VCN preparation, distance/azimuth/elevation pathways, IC/SC readout.}

\section{Signal Model}
\writingbox{Describe the FM call, propagation delay, inverse-square attenuation, head shadow, noise/jitter, and spectral cue insertion.}

\section{Cochlea and Squelch Frontend}
\figuretodo{0.82}{ch3_final_cochlea_output.png}{Final optimised cochlea output.}
\figuretodo{0.82}{ch3_channel_scaling_runtime.png}{Runtime scaling of optimised IIR cochlea against FFT baseline.}
\writingbox{Explain the final IIR + TorchScript LIF + active-window-gated cochlea. Describe why this replaced FFT/IFFT per channel.}

\section{VCN and Latency Calibration}
\figuretodo{0.82}{ch3_cochlea_latency_vector.png}{Cochlea latency vector used to align CD and sensory sweeps.}
\figuretodo{0.82}{ch3_dynamic_schedule.png}{Dynamic threshold and beta schedule for noisy distance cochlea.}
\writingbox{Explain onset detection, cochlear latency correction, dynamic threshold/beta, and why these are needed before pathway processing.}

\chapter{Parallel Processing Pathways}
\section{Distance Pathway: Temporal Coding}
\figuretodo{0.9}{ch4_distance_pipeline.png}{Final distance pathway pipeline.}
\figuretodo{0.9}{ch4_distance_rasters.png}{Frequency-time rasters through the distance pathway.}
\figuretodo{0.72}{ch4_distance_readout_scatter.png}{Distance pathway readout scatter.}
\writingbox{Follow the structure: biological mechanism, neuromorphic abstraction, parameter tuning, isolated validation. Include FM sweep corollary discharge, VCN onset sharpening, DNLL suppression, IC coincidence detection, AC population map, SC readout.}

\section{Azimuth Pathway: Binaural Level and Time}
\figuretodo{0.88}{ch4_azimuth_pipeline.png}{Azimuth pathway pipeline.}
\figuretodo{0.82}{ch4_azimuth_warped_ild.png}{Inverse-sigmoid ILD warping.}
\figuretodo{0.82}{ch4_azimuth_full3d_pm45.png}{Azimuth full-3D test in the tuned $\pm45^\circ$ range.}
\writingbox{Describe ITD and ILD branches, why ILD dominated in the tuned space, the inverse-sigmoid synaptic warp, and the effect of full 3D tests.}

\section{Elevation Pathway: Spectral Shaping}
\figuretodo{0.82}{ch4_elevation_deep_psd.png}{Received PSD before/after deep comb-filter notch.}
\figuretodo{0.82}{ch4_elevation_synaptic_weights.png}{Signal-and-notch synaptic weight matrix.}
\figuretodo{0.72}{ch4_elevation_full3d_scatter.png}{Elevation full-3D prediction scatter.}
\writingbox{Explain comb-filter pinna cue, DCN-style signal-and-notch weighting, first-notch diagnostic, lateral inhibition attempt, dynamic wideband inhibition attempt, and full-3D elevation result.}

\chapter{Cortical Integration and Continuous Attractor Dynamics}
\section{Population Codes and Readout}
\writingbox{Explain why distance, azimuth, and elevation are represented as population codes rather than direct scalar regression.}

\section{Line Attractor Theory}
\figuretodo{0.82}{ch5_cann_chosen_matrices.png}{Chosen finite-line attractor input and recurrence matrices.}
\figuretodo{0.82}{ch5_energy_landscape.png}{Energy landscape visualisation for the line attractor.}
\writingbox{Explain conversion from ring to line attractor, boundary handling, balanced E/I recurrence, Fisher-information-guided input choices, and biological limits on alpha/rate.}

\section{Pathway Integration}
\figuretodo{0.82}{ch5_line_attractor_dynamics.png}{Distance SC line-attractor dynamics.}
\figuretodo{0.7}{ch5_readout_mae_comparison.png}{Comparison of simple COM and line-attractor readouts.}
\writingbox{Discuss where CANN helped, where it did not, and why attractors may be more useful for temporal tracking/noise filtering than for single-frame clean decoding.}

\chapter{System-Level Performance and Discussion}
\section{Distance Results}
\figuretodo{0.75}{ch6_distance_full_test_accuracy.png}{Distance pathway full 3D test result.}
\writingbox{Summarise clean/noisy/full-3D distance performance and compare with old trained models.}

\section{Azimuth Results}
\writingbox{Summarise tuned $\pm45^\circ$ performance, full $\pm90^\circ$ degradation, ILD vs ITD, and biological interpretation.}

\section{Elevation Results}
\figuretodo{0.82}{ch6_elevation_tuning_sweep.png}{Elevation dynamic/lateral inhibition tuning sweep.}
\writingbox{Summarise 1D and full-3D elevation performance. Explain why lateral/dynamic inhibition did not improve clean COM decoding.}

\section{Environmental, Social, and Ethical Considerations}
\writingbox{Discuss low-power sensing, surveillance/dual-use implications, environmental impact of edge sensors, and animal-inspired design.}

\section{Limitations}
\writingbox{Discuss missing full end-to-end integrated 3D test, simplified acoustics, idealised azimuth-gated ear selection, simulated noise, no multi-object clutter handling yet, and limited biological validation.}

\chapter{Conclusions and Future Work}
\writingbox{State the main findings and future steps: integrate all pathways, stress-test under noise/clutter, add object classification, multi-object tracking, and targeted learning.}

\appendix
\chapter{Risk Assessment Retrospective}
\writingbox{One side of A4 maximum. Compare original Safety Office risk assessment with actual encountered risks: mostly computational/data management; low physical hazard; risks around data loss, compute time, and project scope.}

\chapter{Selected Pseudocode or Implementation Details}
\writingbox{Include only novel pseudocode: coordinate event accumulator, dynamic cochlea schedule, signal-and-notch elevation weights, line attractor update. Avoid full code listings.}

\end{document}
